\documentclass[../../main.tex]{subfiles}
\begin{document}
\begin{oframed}
  \subfile{problem}
\end{oframed}

{\bf [解答]}

(イ), (ハ)から $F(x)$ は $n+1$ 次式であり，（ロ）から$F(x)$は$g(x)$で割り切れるので，$a \neq 0, a, b \in \mathbb{R}$ として
\begin{align}
  F(x) = (ax+b) g(x) \label{eq:1}
\end{align}
とおける．両辺を$x$で微分して，(イ)の条件$F'(x)=g(x)$より
\begin{align}
  g(x) = a g(x) + (ax+b) g'(x) \label{eq:2}
\end{align}
を得る．

ここで (ハ)から，$g(x)$は
\begin{align}
  g(x) = x^n + a_{n-2} x^{n-2} + \cdots \label{eq:3}
\end{align}
と展開できるから，両辺$x$で微分して
\begin{align}
  g'(x) = n x^{n-1} + (n-2) a_{n-2} x^{n-3} + \cdots \label{eq:4}
\end{align}
となる．これら\cref{eq:3,eq:4}を\cref{eq:2}に代入して
\begin{align}
  (a-1)\left[x^n + a_{n-2} x^{n-2} + \cdots\right] + (ax+b)\left[n x^{n-1} + (n-3) a_{n-2} x^{n-3} + \cdots\right] \label{eq:5}
\end{align}
である． $x^n, x^{n-1}$ の項の係数を比較して
\begin{align}
  \begin{cases}
    x^n:  0 = (a-1)  + a n  \\
    x^{n-1}:  0 =  b n
  \end{cases} \\
  \therefore
  b = 0, a = \dfrac{1}{1+n} \quad (\because n > 0) \label{eq:6}
\end{align}
を得る．\cref{eq:2}に代入して
\begin{align}
  ng(x) = x g'(x) \label{eq:7}
\end{align}
となる．$y=g(x)$とおいて微分方程式をとくと
\begin{align}
  & \dfrac{n}{x} dx = \dfrac{1}{y} dy \\
  & n \log x + C_1 = \log y \quad (C_1: \text{定数}) \\
  & y = e^{C_1} x^n
\end{align}
(ハ)より $g(x)$ の $x^n$ の係数は 1 だから $e^{C_1} = 1$ となり
\begin{align}
  g(x) = x^n
\end{align}
となる．よって\cref{eq:1}に代入して
\begin{align}
  F(x) = \dfrac{1}{n+1} x^{n+1}
\end{align}
である．以上まとめて
\begin{align}
  F(x) = \dfrac{1}{n+1} x^{n+1},
  \quad
  g(x) = x^n
\end{align}
である．

\newpage

\end{document}
